\section{\tr{Installation}}

\subsection{\tr{Lieferumfang}}
\begin{itemize}
	\item 	\tr{LARUS Vario Display}
	\item 	\tr{Befestigungsschrauben}
	\item 	\tr{1:1 Standard-RJ45-Kabel}
	\item 	\tr{Micro-SD-Karte mit Adapter}
	\item 	\tr{D-Sub-9-Lötstecker und Gehäuse}
	\item 	\tr{1,5 mm Sechskantschlüssel zur Montage der Drehknöpfe}
\end{itemize}

\subsection{\tr{Schnellstartanleitung}}

\tr{In einigen Fällen können die Standard-Einstellungen verwendet werden und das Gerät in einer
	vereinfachten Art und Weise in Betrieb genommen werden. Dann ist es ausreichend, wenn folgende
	Punkte beachtet werden:}
\begin{enumerate}
	\item 	\tr{Bitte demontieren Sie beide Drehknöpfe mit einem 1,5-mm-Sechskantschlüssel (im
			Lieferumfang enthalten).}
	\item 	\tr{Befestigen Sie das LARUS Vario Display mit drei Schrauben in einer 57 mm großen
			Aussparung der Instrumententafel. Es kann auch um 90°, 180° und 270° verdreht montiert
			werden.}
	\item 	\tr{Beide Drehknöpfe montieren}
	\item 	\tr{Verbinden Sie den CAN-Anschluss des LARUS Vario-Displays und der LARUS-Sensorbox mit
			dem im Lieferumfang enthaltenen 1:1-Patchkabel.}
	\item 	\tr{Schalten Sie LARUS ein.}
	\item 	\tr{Überprüfen Sie, ob das Satellitensymbol auf dem Bildschirm gelb oder grün ist und
			die aktuelle Kursrichtung angezeigt wird.}
	\item 	\tr{Wählen Sie eine geeignete Polare für Ihr Segelflugzeug aus oder erstellen Sie eine.}
	\item 	\tr{Ihr LARUS Vario Display ist nun flugbereit.}	
\end{enumerate}

\subsection{\tr{Design und Funktion}}

\tr{Das LARUS Vario Display zeigt die von LARUS gemessenen und berechneten Daten an. LARUS ist ein
fortschrittliches Variometer mit Echtzeit-Windmessfunktion. Es verfügt über modernste Drucksensoren,
eine fortschrittliche IMU und GNSS-Empfänger, um präzise Flugdaten zu erfassen. Die wichtigsten
Merkmale des Displays sind:}
\begin{itemize}
	\item 	\tr{Rundes Display für 57-mm-Standard-Instrumententafelausschnitte-}
	\item 	\tr{Heller und farbenfroher Bildschirm-}
	\item 	\tr{Leichtes, kompaktes Design mit schwarz eloxiertem Aluminiumgehäuse-}
	\item 	\tr{Drehknopf mit zwei Ebenen und Drucktastenfunktion zum Ändern von Einstellungen und
			zum Aufrufen der Menüs 1,5-mm-Sechskantschlüssel.}
\end{itemize}

\tr{Das LARUS Vario Display wird von Prof. Dr. Klaus Schaefer, Maximilian Betz, Winfried Simon,
Peter Simon und dem SteFly-Team entwickelt und kontinuierlich verbessert. Sie können sich gerne 
an der Entwicklung beteiligen und Verbesserungsvorschläge oder Probleme einbringen.}

\href{https://github.com/larus-breeze}{github.com/larus-breeze}


\subsection{\tr{Systemkonfigurationen}}
\label{system-configurations}

\subsubsection{\tr{Konfiguration im Einsitzer}}

\begin{figure}[H]
    \centering
    \includegraphics[width=16cm]{config-singleseater.jpg}
    \caption{\tr{Anschlussplan Einsitzer-Konfiguration}}
    \label{fig:single-seater-configuration}
\end{figure}

\tr{Im Schaubild oben wird der Navigations Rechner an dien LARUS-Sensorbox angeschlossen. Oftmals
ist es von der Verdrahtung einfacher, den Navigations Rechner an das LARUS Vario Display
anzuschließen. Dies ist genauso gut möglich.}

\subsubsection{\tr{Konfiguration im Doppelsitzer}}

\begin{figure}[H]
    \centering
    \includegraphics[width=16cm]{config-doubleseater.jpg}
    \caption{\tr{Anschlussplan Doppelsitzer-Konfiguration}}
    \label{fig:double-seater-configuration}
\end{figure}

\tr{Im Doppelsitzer werden die Navigationsrecher typischerweise an die LARUS Vario Displays
angeschlossen. Es ist aber genauso möglich, die Navigationsrechner an die LARUS-Sensorbox
anzuschließen. Die Daten werden zwischen den LARUS Komponenten ausgetauscht, so dass keine
Informationen verloren gehen.}

\subsection{\tr{Steckverbinder und Verkabelung}}

\subsubsection{\tr{Steckverbinder auf der Rückseite des Gerätes}}

\begin{figure}[H]
    \includegraphics[width=8cm]{connectors-overview.jpg}
    \caption{\tr{Ansicht von der Rückseite des Gerätes}}
    \label{fig:connectors-overview}
\end{figure}

\tr{Auf der Rückseite des LARUS Vario Displays befindesn sich die Anschlüsse für CAN, RS232 und die
Ein-/Ausgänge. Des weiteren ist dort der Schlitz zum Einführen der SD-Karte vorhanden.}

\subsubsection{\tr{CAN- und RS232-Anschlüsse}}

\begin{figure}[H]
    \includegraphics[width=8cm]{connectors-rj45.pdf}
    \caption{\tr{Die Rj45 Anschlüsse im Detail}}
    \label{fig:connectors-rj45}
\end{figure}

\begin{table}[H]
	\begin{tabular}{C{0.1\linewidth} L{0.40\linewidth} L{0.40\linewidth}}
		\toprule
		\textbf{Pin} & \textbf{CAN} & \textbf{RS232} \\
		\midrule
		1 & GND (internally connected) & GND (internally connected) \\
		2 & GND (internally connected) & GND (internally connected) \\
		3 & NC 						   & RS232-1-RX \\
		4 & CAN Low                    & RS232-1-TX \\
		5 & CAN High                   & NC \\
		6 & NC                         & NC \\
		7 & VCC [9-28V DC]             & VCC [9-28V DC] \\
		8 & VCC [9-28V DC]             & VCC [9-28V DC] \\
		\bottomrule
	\end{tabular}
	\caption{\tr{Pinbelegung CAN und RS232 RJ45}}
\end{table}

\subsubsection{GPIO / D-SUB 9}
\tr{Über den D-Sub-Stecker können mehrere weitere Schalter, Sensoren und Geräte angeschlossen
werden. Bitte beachten Sie, dass zum Zeitpunkt der Veröffentlichung dieses Handbuchs noch nicht alle
Funktionen implementiert sind. Die folgende Skizze zeigt die Ansicht in den Stecker des LARUS Vario
Displays.}

\begin{figure}[H]
    \includegraphics[width=6cm]{connector-dsub9.png}
    \caption{\tr{DSUB-9 Anschluss}}
    \label{fig:connectors-dsub9}
\end{figure}

\begin{table}[H]
	\begin{tabular}{C{0.1\linewidth} L{0.40\linewidth} L{0.20\linewidth}}
		\toprule
		\textbf{Pin} & \textbf{CAN} & \textbf{RS232} \\
		\midrule
		1 & GND  					& (Ground) \\
		2 & DI3 - Gear 				& Input \\
		3 & DI1 - Water Ballast		& Input \\
		4 & DO2						& Output \\
		5 & GND						& (Ground) \\
		6 & DI4 - Speed Breakes 	& Input \\
		7 & DI2 - Spped to fly		& Input \\
		8 & GND						& (Ground) \\
		9 & DO1 - Canopy Flasher 	& Output \\
		\bottomrule
	\end{tabular}
	\caption{Pin Asssignment GPIO / D-SUB 9}
\end{table}

\tr{Zur leichteren Identifizierung der Pins sind diese Nummern auch in die Buchse (im Lieferumfang
enthalten) eingeprägt. Nach der Verkabelung müssen die Einstellungen im LARUS Vario Display
vorgenommen werden.}

\subsubsection{\tr{Audio}} \tr{Eine Audiobuchse ist für den Anschluss eines Lautsprechers mit einem
3,5-mm-Klinkenstecker verfügbar. Der Innenwiderstand des Lautsprechers muss zwischen 4 und 8 Ohm
liegen (max. output of 3 W bei 8 Ohm).}

\picnote
	{pictograph-yellow-warning.pdf}
	{\tr{Ein einzelner Lautsprecher darf nicht an mehr als ein Gerät angeschlossen werden.}}

\subsubsection{\tr{SD-Karte}}
\tr{Das Gerät verfügt über einen SD-Kartensteckplatz für Firmware-Updates.}

\picnote
	{pictograph-yellow-warning.pdf} {\tr{Da SD-Kartenerweiterungen das LARUS Vario Display
	beschädigen können, übernehmen wir keine Haftung für Schäden, die durch deren Verwendung
	entstehen.}}

\subsubsection{\tr{Verkabelung}}

\tr{Wenn eine serielle Verbindung (RS232) zwischen SteFly NAV und LARUS Vario Display hergestellt
wird, werden beispielsweise die MC-Werteinstellungen zwischen den Geräten synchronisiert/übertragen.
Darüber hinaus verarbeitet das LARUS Vario Display Eingaben von einem Speed-to-Fly-/Vario-Schalter,
der direkt mit dem SteFly-Fernsteuerungshebel verbunden ist. Die folgenden Schritte sind
erforderlich:}
\begin{itemize}
	\item 	\tr{Geräte mit Kabel verbinden:}
	\begin{itemize}
		\item 	\tr{Der CAN-Anschluss der LARUS Sensor Box muss über 1:1-Patchkabel mit dem
				CAN-Anschluss des LARUS VARIO Displays verbunden werden (bei
				Doppelsitzerkonfigurationen unter Verwendung der Kabelbox).}
		\item 	\tr{Die Geräte, auf denen XCSoar läuft, können über ein gekreuztes Patchkabel RX/TX
				zwischen dem RS232-Anschluss des Vario-Displays und dem ttySx-Anschluss des SteFly
				NAV verbunden werden.}
	\end{itemize}
	\item 	\tr{Schließen Sie optional Schalter für die Erkennung des Ablassvorgangs beim
			Wasserballasts und die Fahrwerkswarnung an den entsprechenden DSUB-9-Pin und einen
			beliebigen GND-Pin an. Für die Fahrwerkswarnung haben Sie die Möglichkeit, zwei
			separate, direkt angeschlossene Schalter für das Fahrwerk und die Bremsklappen zu
			verwenden (empfohlen) oder ein Kabel mit den beiden Schaltern in Reihe zum
			DSUB-9-Anschluss.}
	\item 	\tr{Der SteFly-Fernbedienungsstick kann auf zwei verschiedene Arten angeschlossen
			werden:}
	\begin{itemize}
		\item 	\tr{Abhängig von der Ausstattung: mit einem separaten Kabel zum DSUB-9-Pin 7}
		\item 	\tr{Verwendung des USB-Kabels (immer verfügbar)}
	\end{itemize}
	\item	\tr{Bitte passen Sie die Einstellungen in XCSoar (Version 7.44 oder höher) / OpenSoar
			(7.43 oder höher) an:}
	\begin{itemize}
		\item 	\tr{Wenn das Vario-Display direkt an das SteFly NAV angeschlossen ist,
		wählen Sie bitte den entsprechenden ttyS-Port aus: Konfiguration -> Geräte -> wählen
		Sie eine freie Leitung aus und geben Sie den Port ein, Baudrate 38400, Treiber
		Larus, Synchronisation (Option) mit Gerät aktivieren.}
	\end{itemize}
	\item 	\tr{Bitte passen Sie die Einstellungen im LARUS Vario Display für die angeschlossenen
			optionalen Schalter an.}
	\begin{itemize}
		\item 	\tr{Warnung bezüglich der Ausrüstung}
		\item 	\tr{Wasserballast}
		\item 	\tr{StF-Schalterkonfiguration (falls nicht automatisch geschaltet)}
		\begin{itemize}
			\item 	\tr{Eingangs-Pin: Um einen externen Schalter (am Steuerknüppel, Klappenhebel
					usw.) zu verwenden, stellen Sie die Vario-Steuerung auf „Eingangs-Pin“ und
					wählen Sie die richtige Einstellung in der StF-Pin-Konfiguration. (active: when
					open/closed)}
			\item 	\tr{NMEA für SteFly-Fernsteuerungsstick. Zusätzlich zu den Einstellungen im
					Vario müssen Sie eine Ereignisdefinitionsdatei in XCSoar/OpenSoar installieren
					(Konfiguration / System / Ansicht / Sprache, Eingabe / Ereignisse – Expert
					aktivieren, auf Ereignisse klicken – Herunterladen –
					GLB-XCI-xcremote-XCNAV.xci). Verlassen Sie die XCSoar-Konfiguration und starten
					Sie XCSoar neu.}
		\end{itemize}
	\end{itemize}
\end{itemize}

\subsection{\tr{CAN-Terminierung}}

\tr{Das LARUS Vario Display und die LARUS Sensor Box sind über den CAN Bus miteinander verbunden.
CAN-Bus-Netzwerke erfordern Abschlusswiderstände an jedem Ende des Netzwerks. Daher verfügen alle
Geräte über einen integrierten Schalter zur Aktivierung des Widerstands:}

\begin{figure}[H]
    \centering
    \includegraphics[width=16cm]{can-termination.jpg}
    \caption{\tr{Schalter an den Geräten zur Terminierung des CAN Bus}}
    \label{fig:can-termination}
\end{figure}

\begin{table}[H]
	\begin{tabular}{L{0.25\linewidth} C{0.156\linewidth} C{0.156\linewidth} C{0.156\linewidth} C{0.156\linewidth}}
		\toprule
		\textbf{Description} & \textbf{Display front} & \textbf{Larus Box} & \textbf{CAN Splitter} & \textbf{Display rear} \\
		\midrule
		\tr{Einsitzer} 							 & on  & on  & -   & -   \\
		\tr{Zweisitzer, Larus-Box in der Frontblende} & off & on  & off & on  \\
		\tr{Zweisitzer, Larus-Box in der Rückwand}  & on  & on  & off & off \\
		\bottomrule
	\end{tabular}
	\caption{\tr{CAN-Terminierung Ein-/Zweisitzer}}
\end{table}

\tr{Bitte beachten Sie: Alle LARUS-Sensorboxen, die vor März 2025 ausgeliefert wurden, verfügen über
keinen CAN-Abschlussschalter. Die CAN-Abschlusswiderstände sind standardmäßig immer aktiviert.}

\subsection{\tr{Externe Sicherung}}

\picnote
	{pictograph-yellow-warning.pdf} 
	{\tr{In der Regel wird das LARUS Vario Display über ein
	Patchkabel zwischen den CAN-Anschlüssen von LARUS mit Strom versorgt. LARUS muss durch eine
	externe Sicherung (500 mA bis max. 3 A) geschützt werden, wie es bei allen elektrischen Geräten
	in der Luftfahrt üblich ist. Wenn LARUS seine Energie von einem anderen Hauptinstrument bezieht
	(z. B. SteFly NAV über D-Sub-Stecker), stellen Sie bitte sicher, dass das Hauptinstrument durch
	eine externe Sicherung entsprechend geschützt ist.}}

\subsection{\tr{Installationsort}}

\tr{Das folgende Bild zeigt eine typische Einbausituation des LARUS Vario Displays im
Instrumentenbrett eines Segelflugzeugs.}

\begin{figure}[H]
    \includegraphics[width=8cm]{panel.jpg}
    \caption{\tr{Einbausituation im Instrumentenpanel}}
    \label{fig:panel}
\end{figure}

\tr{Das Display passt in eine Standardaussparung von 57 mm und wird mit 3 M3-Schrauben befestigt.}

\begin{figure}[H]
    \includegraphics[width=8cm]{knobs-and-screws.jpg}
    \caption{\tr{Befestigungselemente}}
    \label{fig:knobs-and-screws}
\end{figure}

\tr{Für die Installation ist es erforderlich, die beiden Drehknöpfe mit einem
1,5-mm-Sechskantschlüssel zu entfernen.}

\subsection{\tr{Installationsorientierung}}

\picnote
	{pictograph-yellow-warning.pdf} {\tr{Wenn Sie beabsichtigen, ein vorhandenes Gerät durch das
	Lars Vario Display zu ersetzen, überprüfen Sie bitte die gewünschte Einbaulage, bevor Sie das
	7,3-mm-Loch für den Drehgeber bohren, da das Gehäuse des Displays leicht asymmetrisch ist.}}

\tr{Das Display kann in den Ausrichtungen 0° / 90° / 180° / 270° montiert werden. Nach der Montage
des Displays muss dessen Ausrichtung im Menü „Display Rotation“ ggfs. angepasst werden.}

\subsection{\tr{Erstinbetriebnahme und Funktionstest}}

\tr{Für die Inbetriebnahme befolgen Sie bitte die folgenden Schritte:}
\begin{enumerate}
	\item 	\tr{Bitte überprüfen Sie, ob das LARUS Vario Display gemäß den Zeichnungen in Kapitel }
			\nameref{system-configurations} \tr{ an die LARUS Sensor Box oder den CAN Splitter
			angeschlossen ist.}
	\item 	\tr{LARUS einschalten}
	\item	\tr{Bitte überprüfen Sie, ob LARUS Vario Displays startet und ein gelbes oder grünes
			Satellitenpiktogramm angezeigt wird. Die Vario-Zeiger sollten leichte Bewegungen um die
			Nullposition ausführen.}
\end{enumerate}

\subsection{\tr{Wartung}} \tr{Das gesamte System enthält keine zu wartenden Teile. Um
Garantieleistungen in Anspruch zu nehmen, wenden Sie sich bitte direkt an SteFly.}

\picnote
	{pictograph-yellow-warning.pdf}
	{\tr{Das Öffnen des Gehäuses des LARUS Vario Displays führt zum Erlöschen der Garantie.}}

\subsection{\tr{Fehlerbehebung}}
\begin{table}[H]
	\begin{tabular}{L{0.29\linewidth} L{0.29\linewidth} L{0.29\linewidth}}
		\toprule
		\textbf{Problem} & \textbf{Possible Causes} & \textbf{Solution} \\
		\midrule
		\tr{Das LARUS Vario Display startet, jedoch ist das Satellitenpiktogramm rot und die Vario-Zeiger sind fixiert.}
		& \tr{Anschluss an den LARUS-CAN-Port mit einem gekreuzten Rx/Tx-Patchkabel anstelle eines standardmäßigen 1:1-Patchkabels.}
		& \tr{Bitte ersetzen Sie das Patchkabel und verwenden Sie das im Lieferumfang enthaltene Kabel.} \\
		\tr{Das LARUS Vario Display startet, jedoch ist das Satellitenpiktogramm rot und die Vario-Zeiger sind fixiert.}
		& \tr{Anschluss an LARUS über RS232-Port}
		& \tr{Bitte verbinden Sie die CAN-Anschlüsse.} \\
		\bottomrule
	\end{tabular}
	\caption{\tr{Fehlerbehebung}}
\end{table}

\subsection{\tr{Firmware-Aktualisierung}}

\tr{Das LARUS-Team verbessert die Software kontinuierlich und veröffentlicht Updates, beispielsweise
auf der LARUS Vario Display-Webseite. Um die Firmware zu aktualisieren, speichern Sie bitte die neue
.bin-Datei auf der im Lieferumfang enthaltenen SD-Karte und stecken Sie diese in den
SD-Kartensteckplatz des LARUS Vario Display. Das Gerät installiert nur kompatible
Firmware-Versionen, die höher sind als die installierte Version. Nach dem Einschalten startet das
Gerät in weniger als einer Sekunde. Wenn eine kompatible Firmware auf der SD-Karte erkannt wird,
bleibt das Display etwa 3-5 Sekunden lang schwarz, bevor die Meldung „Firmware-Upgrade wird
durchgeführt“ angezeigt wird. Anschließend wird die geladene Firmware-Version angezeigt.

Sollte das Update nach dem Einschalten des Geräts mit einer eingelegten SD-Karte, die eine neue
Firmware-Version enthält, nicht gestartet werden, verwenden Sie bitte eine andere SD-Karte.
Anforderungen an die SD-Karte: min. 4 GB SDHC, FAT32-Format (nicht GPT).}
